\subsection{Full-AA quantifiers and smooth-to-set rounding}
\label{ian:hard-sector-sets-asymptotics}

Status: proved, with no uniform-AA hypothesis (full AA appears only as the property being characterized): two reformulations of Theorem~\ref{ian:hard-sector-aa-equivalence}, one with unspecified polynomial weights and one resource-sensitive, for smooth and for set sectors; and, as side results independent of extraction, a rounding criterion under an atom-size bound and an atomic obstruction without one.

Throughout, the ambient law is \(\mu_n\) and \(\mathcal Q_{n,T}\) is as
before \eqref{ian:hard-sector-aa-definition}. A
\((w,t,a,a/4)\)-hard sector of \(Q\) is a sector of \(\mu_n\) of weight
\(w\), smooth or set according to the interpretation fixed, whose law
\(\nu\) satisfies \(\operatorname{Hard}(Q,\nu;t,a,a/4)\).

\begin{remark}[An equivalent formulation using unspecified polynomials]
\label{ian:hard-sector-polynomial-quantifiers}
For either smooth sectors or set sectors, failure of full AA is
equivalent to the following assertion:
\[
 \begin{gathered}
 \exists K,p,q\ge1\quad\forall A,B\ge1\quad
 \exists n,T\ge1,\ r>8,\ Q\in\mathcal Q_{n,T},\
 t\in\N,\ w\in(0,1],\ a\in(0,1/4]:\\
 t>\left\lceil A(T+r)^B\right\rceil,\qquad
 w^{-1}\le Kr^p,\qquad a^{-1}\le Kr^q,
 \end{gathered}
 \tag{P}\label{ian:hard-sector-polynomial-property}
\]
together with the existence of a \((w,t,a,a/4)\)-hard sector of \(Q\) in
the chosen interpretation.

The canonical witnesses in
Theorem~\ref{ian:hard-sector-aa-equivalence}, items (2) and (3),
establish \eqref{ian:hard-sector-polynomial-property} with
\(K=64,p=3,q=2\), since \(1/(16r^2)=a/4\) for \(a=1/(4r^2)\).

Conversely, suppose \eqref{ian:hard-sector-polynomial-property} holds
and full AA has constants \(A_0,B_0\). Let \(s=p+q\) and choose
\[
 A=A_0(1+8K^2/3)^{B_0}+1,\qquad B=sB_0.
\]
For its witness, which is a smooth sector in either interpretation,
\(8/(3wa)\le(8K^2/3)r^s\). Since \(T+r\ge1\) and \(s\ge1\), one has
\(T\le(T+r)^s\) and \(r^s\le(T+r)^s\), so
\(T+cr^s\le(1+c)(T+r)^s\) for \(c\ge0\); with \(c=8K^2/3\),
\(T+8/(3wa)\le(1+8K^2/3)(T+r)^s\). Lemma~\ref{ian:hard-sector-aa-exclusion}
then gives
\(t<\lceil A_0(1+8K^2/3)^{B_0}(T+r)^{sB_0}\rceil\le\lceil A(T+r)^B\rceil\),
contradicting \(t>\lceil A(T+r)^B\rceil\). The polynomial bounds must be
fixed before the budget quantifiers.
\end{remark}

\begin{corollary}[The resource-sensitive sector formulation]
\label{ian:hard-sector-resource-equivalence}
Fix either the smooth or the set interpretation of sectors. For a family
\(\symbf{\mu}\), let \(\mathsf R(\symbf{\mu})\) mean
\[
 \begin{gathered}
 \forall b>0\quad\exists A_b,B_b\ge1\quad
 \forall n,T\ge1\quad\forall Q\in\mathcal Q_{n,T}\quad
 \forall a\in(0,1/4]\quad\forall w\in(0,1]\quad
 \forall t\in\N:\\
 t\ge\left\lceil A_b(2+T+a^{-1})^{B_b}\right\rceil
 \quad\Longrightarrow\quad
 \text{every }(w,t,a,a/4)\text{-hard sector has }
 w<(2+T+a^{-1})^{-b}.
 \end{gathered}
\]
Then full AA holds if and only if
\(\mathsf R(\symbf{\mu})\) holds, in either interpretation.
\end{corollary}

\begin{proof}
Assume full AA with constants \(A_0,B_0\), fix \(b>0\), and set
\[
 A_b=A_0\,4^{B_0}+1,\qquad B_b=(b+1)B_0.
\]
Suppose a \((w,t,a,a/4)\)-hard sector has \(w\ge z^{-b}\), where
\(z=2+T+a^{-1}\), and \(t\ge\lceil A_bz^{B_b}\rceil\). Since \(z\ge T\),
\(z\ge a^{-1}\) and \(z>1\),
\[
 T+\frac8{3wa}
 \le T+\frac83a^{-1}z^b
 \le z+\frac83z^{b+1}
 \le\frac{11}{3}z^{b+1}<4z^{b+1},
\]
so Lemma~\ref{ian:hard-sector-aa-exclusion} gives
\(t<\lceil A_0\,4^{B_0}z^{(b+1)B_0}\rceil\le\lceil A_bz^{B_b}\rceil\le t\),
a contradiction. Hence \(\mathsf R\) holds.

For the converse, suppose full AA fails; we show that \(\mathsf R\) fails
already for \(b=3\). Fix arbitrary proposed constants \(A,B\ge1\), and
apply Theorem~\ref{ian:hard-sector-aa-equivalence}, item (2) or (3)
according to the interpretation, with budget constants
\[
 A'=A7^B,\qquad B'=2B.
\]
It supplies a \((w,t,a,a/4)\)-hard sector with
\[
 a=\frac1{4r^2},\qquad w=\frac1{64r^3},\qquad
 t=\left\lceil A7^B(T+r)^{2B}\right\rceil+1.
\]
For \(z=2+T+a^{-1}=2+T+4r^2\), since \(T\ge1\) gives \(7T^2\ge T+2\),
\[
 z\le7T^2+7r^2\le7(T+r)^2,
 \qquad\text{hence}\qquad
 t>\left\lceil A7^B(T+r)^{2B}\right\rceil\ge\left\lceil Az^B\right\rceil.
\]
Moreover,
\[
 z^{-3}\le(4r^2)^{-3}
 =\frac1{64r^6}\le\frac1{64r^3}=w.
\]
This sector violates the conclusion of \(\mathsf R\) for \(b=3\),
\(A_3=A\), \(B_3=B\). Since \(A,B\) were arbitrary, \(\mathsf R\) fails.
\end{proof}

The last two results concern rounding a supplied smooth sector to a set
sector, which can require an atom-size hypothesis. Extraction
(Corollary~\ref{ian:hard-sector-explicit}) produces set sectors directly
and uses neither.

\begin{lemma}[Rounding a smooth sector under an atom-size bound]
\label{ian:hard-sector-rounding}
Let \(n,t\ge1\), \(Q:\Omega_n\to[0,1]\), and let \(\mu,\nu\) be
probability laws with \(\nu\le\mu/w\), \(0<w\le1\). Put
\[
 N=n^{2^t-1}2^{2^t},\qquad J=3+2N,\qquad
 \alpha=\max_x\mu(x),
\]
and let \(0<\delta\le1/2\) with
\(\alpha\le\delta^2w^2/\log(4J)\).
There is \(H\subseteq\Omega_n\) with \(\mu(H)\ge(1-\delta)w\ge w/2\) such
that, for \(\lambda=\mu(\cdot\mid H)\),
\[
 \Var_\lambda Q\ge\Var_\nu Q-12\delta,\qquad
 |\Cov_\lambda(Q,F)|\le|\Cov_\nu(Q,F)|+12\delta\quad(F\in\mathcal F_t).
\]
Consequently \(\operatorname{Hard}(Q,\nu;t,a,\gamma)\) implies
\(\operatorname{Hard}(Q,\lambda;t,a-12\delta,\gamma+12\delta)\)
whenever \(0<\gamma+12\delta<a-12\delta\).
\end{lemma}

\begin{proof}
Independently select each point \(x\) with \(\mu(x)>0\) with probability
\(w\nu(x)/\mu(x)\in[0,1]\), and let \(H\) be the selected set; points
with \(\mu(x)=0\) have \(\nu(x)=0\). Let \(\mathcal L\) consist of
\(1,Q,Q^2\) and \(F,QF\) for the at most \(N\) Boolean-valued depth-\(t\)
tests \(F\) (Lemma~\ref{ian:hard-sector-tree-closure}): at most \(J\)
functions, each bounded by one in absolute value. For
\(f\in\mathcal L\) put \(Z_f=\sum_x\mu(x)\mathbf1_H(x)f(x)\), so
\(\E Z_f=w\E_\nu f\). The summands are independent with ranges of length
at most \(\mu(x)\), and \(\sum_x\mu(x)^2\le\alpha\), so Hoeffding's
inequality gives
\[
 \Pr\{|Z_f-w\E_\nu f|>\delta w\}\le2e^{-2\delta^2w^2/\alpha}\le2(4J)^{-2}.
\]
A union bound over \(\mathcal L\) leaves probability at least
\(1-1/(8J)>0\) that every deviation is at most \(\delta w\); fix such a
realization. Then \(s=\mu(H)=Z_1\in[(1-\delta)w,(1+\delta)w]\), and for
every \(f\in\mathcal L\),
\[
 |\E_\lambda f-\E_\nu f|
 =\left|\frac{Z_f}{s}-\E_\nu f\right|
 \le\frac{|Z_f-w\E_\nu f|+|w-s|\,|\E_\nu f|}{s}
 \le\frac{2\delta}{1-\delta}\le4\delta.
\]
With \(m_\lambda=\E_\lambda Q\) and \(m_\nu=\E_\nu Q\in[0,1]\),
\(|\Var_\lambda Q-\Var_\nu Q|\le|\E_\lambda Q^2-\E_\nu Q^2|+2|m_\lambda-m_\nu|\le12\delta\),
and for every Boolean-valued test \(F\),
\(|\Cov_\lambda(Q,F)-\Cov_\nu(Q,F)|\le|\E_\lambda QF-\E_\nu QF|+|m_\lambda-m_\nu|+|\E_\lambda F-\E_\nu F|\le12\delta\).
For general \(F\in\mathcal F_t\) the covariance difference is linear in
the leaf outputs of \(F\), so as in
Lemma~\ref{ian:hard-sector-tree-closure} its absolute value is maximized
at a Boolean assignment; hence the bound holds for all
\(F\in\mathcal F_t\).
\end{proof}

\begin{proposition}[An atomic obstruction to general smooth-to-set rounding]
\label{ian:hard-sector-rounding-counterexample}
There exist \(\mu,Q,\nu\) on \(\Omega_2\) such that
\[
 \nu\le2\mu,\qquad
 \operatorname{Hard}(Q,\nu;1,1/4,1/64),
\]
but no positive-mass conditional-set law has both positive variance and
depth-one covariance bounded by \(1/64\).
\end{proposition}

\begin{proof}
Index the four points by
\(x^{(1)}=(1,1)\), \(x^{(2)}=(1,-1)\), \(x^{(3)}=(-1,1)\),
\(x^{(4)}=(-1,-1)\), and set
\[
 (\mu_1,\mu_2,\mu_3,\mu_4)
 =\left(\frac12,\frac14,\frac18,\frac18\right),
 \qquad Q(x)=\frac{1+x_1x_2}{2},
\]
with \(\nu\) uniform on the four points.
Then \(\nu\le2\mu\), \(\Var_\nu Q=1/4\), and
\(\Cov_\nu(Q,x_1)=\Cov_\nu(Q,x_2)=0\). Every depth-one test is an
affine function of one coordinate, so the smooth hardness follows.

A set has positive conditional variance precisely when it contains
both a point with \(Q=1\) (index \(1\) or \(4\)) and a point with
\(Q=0\) (index \(2\) or \(3\)). For each of the nine such subsets, the
maximum absolute coordinate covariance under its conditional law,
computed from
\[
 \Cov_\lambda(Q,x_i)
 =p(1-p)\left(
 \E_\lambda[x_i\mid Q=1]-\E_\lambda[x_i\mid Q=0]\right),
 \qquad p=\lambda(Q=1),
\]
is
\[
\begin{array}{c|c}
 H\text{, specified by point indices}
 &\max_{i=1,2}|\Cov_{\mu(\cdot\mid H)}(Q,x_i)|\\ \hline
 12&4/9\\
 13&8/25\\
 24&4/9\\
 34&1/2\\
 123&16/49\\
 124&16/49\\
 134&2/9\\
 234&1/4\\
 1234&7/32
\end{array}
\]
Every entry is at least \(7/32>1/64\). This does not obstruct
extraction from approximation failure: this \(Q\) is computed exactly
with two classical queries.
\end{proof}
